\subsection{循环}

% Page 19
\begin{frame}{循环：行动、观察与决策}
\begin{center}
\resizebox{.91\textwidth}{!}{\begin{tikzpicture}[b/.style={draw=blue,fill=blue!5,rounded corners,text width=2.9cm,minimum height=.8cm,align=center,font=\small},a/.style={-{Stealth},thick}]
\node[b](p)at(0,1.2){规划下一步};
\node[b](t)at(4,1.2){调用获准工具};
\node[b](o)at(8,1.2){观察 + 决策};
\node[b](r)at(2,-0.8){在预算内继续};
\node[b](e)at(8,-0.8){返回或暂停};
\draw[a](p)--(t);\draw[a](t)--(o);\draw[a](o.south)--(r.east);\draw[a](r.west)-| (p.south);\draw[a](o)--(e);\end{tikzpicture}
}
\end{center}
\begin{itemize}
\item 工具调用失败后，任务仍未完成。
\item 根据观察结果选择下一步行动。
\item 对暂时性错误进行有限次数的重试；无效证据不得推动案例进入下一步。
\end{itemize}
\end{frame}

% Page 20
\begin{frame}{循环：生成、检查与审查}
\setlength{\topsep}{0pt}
\begin{center}
\resizebox{.96\textwidth}{!}{\begin{tikzpicture}[
  box/.style={draw=blue,fill=blue!5,rounded corners=2pt,text width=2.65cm,minimum height=1.4cm,align=center,font=\footnotesize,inner sep=4pt},
  flow/.style={-{Stealth[length=4pt]},line width=.7pt,draw=black!60},
  lab/.style={font=\scriptsize,text=black!60,fill=white,inner sep=2pt}]
\node[box] (g) at (1.6,0) {\textbf{Generator}\\候选策略 + 测试};
\node[box,draw=forestgreen,fill=forestgreen!5] (t) at (5.05,0) {\textbf{代码检查关卡}\\经由 Orchestrator\\必要检查};
\node[box,draw=darkred,fill=red!4] (c) at (8.5,0) {\textbf{新上下文的 Critic}\\核查证据};
\node[box,draw=forestgreen,fill=forestgreen!5] (o) at (11.95,0) {\textbf{Orchestrator}\\修订、验收\\或停止};
\draw[flow] (g)--(t);\draw[flow] (t)--(c);\draw[flow] (c)--(o);
\draw[flow] (o.south)--(11.95,-1.35)--node[below,lab]{修复任务包：问题项 + 所需证据}(1.6,-1.35)--(g.south);
\node[box,text width=3.3cm,minimum height=1.0cm] (claim) at (2.1,-2.75) {\textbf{Generator 的声明}\\\lq\lq 已计入成本\rq\rq};
\node[box,text width=3.3cm,minimum height=1.0cm] (inspect) at (6.8,-2.75) {\textbf{Critic 的核查}\\保存的交易记录 + 回测};
\node[box,text width=3.3cm,minimum height=1.0cm,draw=darkred,fill=red!4] (finding) at (11.5,-2.75) {\textbf{问题项}\\未扣除成本};
\draw[flow] (claim)--(inspect);\draw[flow] (inspect)--(finding);
\node[font=\scriptsize,text=black!55] at (6.8,-3.6) {审查示例};
\end{tikzpicture}
}
\end{center}
\begin{itemize}
\item 向 Critic 提供已保存的交易记录和回测输出，供其核查。
\item 每轮启用新上下文的 Critic 可减轻对先前回答的锚定；但不同 Agent 仍可能犯同样的错误。
\item 修订过程中保持任务授权与验收规则不变。
\end{itemize}
\end{frame}

% Page 21
\begin{frame}[t]{循环：问题项与修复核验\handsup}
\setlength{\topsep}{0pt}
\begin{center}
\resizebox{.94\textwidth}{!}{\begin{tikzpicture}[
  card/.style={draw=blue,fill=blue!4,rounded corners=3pt,align=left,font=\small,inner sep=7pt}]
\node[card] at (3.1,0) {
  \parbox[t][3.55cm][t]{5.65cm}{\raggedright
  \textbf{时序问题}\par\smallskip
  信号使用第 $t$ 日的收盘价。\par\smallskip
  回测却在同一收盘时点交易。\par\smallskip
  交易必须发生在信号可用之后。
  }
};
\node[card,draw=darkred,fill=red!3] at (9.9,0) {
  \parbox[t][3.55cm][t]{5.65cm}{\raggedright
  \textbf{问题项台账}\par\smallskip
  第 1 轮：记录 \texttt{pit-001}，并链接至受影响的输出。\par\smallskip
  第 2 轮遗漏该项：问题仍保持未解决状态。\par\smallskip
  只有 Critic 核验修复后的输出，才能关闭该问题项。
  }
};
\end{tikzpicture}
}
\end{center}
\begin{center}\illustrative\end{center}
\begin{itemize}
\item 问题：对已保存文件做什么检查，才能证明时序问题已经修复？
\item 下一次请求须包含所有尚未解决的问题项。
\end{itemize}
\end{frame}

% Page 22
\begin{frame}{循环：验收与完成}
\begin{center}
\resizebox{.95\textwidth}{!}{\begin{tikzpicture}[
  box/.style={draw=blue,fill=blue!5,rounded corners=2pt,text width=2.7cm,minimum height=2.05cm,align=center,font=\small,inner sep=5pt},
  flow/.style={-{Stealth[length=5pt]},thick,draw=black!60}]
\node[box] (g) at (1.5,0) {\textbf{数值检查关卡}\\全部必要检查通过};
\node[box] (c) at (5,0) {\textbf{完整审查}\\全部 Critic 通过\\且无未解决问题项};
\node[box] (r) at (8.5,0) {\textbf{同一候选策略}\\连续两轮审查通过};
\node[box,draw=forestgreen,fill=forestgreen!6] (a) at (12,0) {\textbf{ACCEPTED}\\研究任务结束};
\draw[flow] (g)--(c);\draw[flow] (c)--(r);\draw[flow] (r)--(a);
\node[font=\small,text=black!60] at (6.75,-1.65) {候选策略发生变化后，连续通过审查的轮数重新计算。};
\end{tikzpicture}
}
\end{center}
\begin{itemize}
\item QT 要求同一候选策略连续两轮通过全部审查；更换候选策略后，连续通过轮数重新计算。
\item 这套课堂验收规则既不能证明策略盈利，也不构成交易授权。
\end{itemize}
\end{frame}

% Page 23
\begin{frame}{循环：预算、路由与停止}
\begin{center}
\small
\begin{tabular}{@{}P{3.5cm}P{9.0cm}@{}}
\toprule
\textbf{路由 / 结果} & \textbf{条件} \\
\midrule
\texttt{CONTINUE}\newline（非终止路由） & 仍可修复或确认，仍有进展且预算未耗尽。 \\
\rowcolor{gray!10}
\texttt{REJECTED} & 出现致命缺陷或 Critic 给出否决结论，终止当前候选方案。 \\
\texttt{BLOCKED} & 限定次数的修复后，必需提交仍无效或必要检查仍未通过。 \\
\rowcolor{gray!10}
\texttt{EXHAUSTED} & 验收通过前已达到预设预算上限。 \\
\texttt{STALLED} & 多轮未通过且没有新进展。 \\
\rowcolor{gray!10}
\texttt{USER\_STOPPED} & 用户明确要求结束运行。 \\
\bottomrule
\end{tabular}

\end{center}
\begin{itemize}
\item 开始前设定轮数、尝试次数、工具调用次数和用时上限。
\item 记录停止原因。预算耗尽绝不能将失败结果变成验收通过。
\end{itemize}
\end{frame}

% Page 24
\begin{frame}{循环：恢复与安全重试}
\begin{center}
\small
\begin{tabular}{@{}P{4.0cm}P{8.5cm}@{}}
\toprule
\textbf{观察到的问题} & \textbf{下一步核查或行动} \\
\midrule
超时 & 有限次数重试前，先检查工作是否已完成。 \\
\rowcolor{gray!10}
格式或覆盖范围错误 & 将具体错误交回生成该报告的角色。 \\
已封存证据发生变化 & 阻止流程推进并调查。 \\
\rowcolor{gray!10}
外部操作是否完成不确定 & 重复操作前，先核对实际影响。 \\
\bottomrule
\end{tabular}

\end{center}
\begin{itemize}
\item 仅恢复已核验的待办工作；已完成的案例直接返回保存的结果。
\item 重试前先核对结果不确定的外部操作，避免重复下单或付款。
\end{itemize}
\end{frame}

\subsection{图}

% Page 25
\begin{frame}{图：依赖关系与交接}
\begin{center}
\resizebox{.94\textwidth}{!}{\begin{tikzpicture}[
  box/.style={draw=blue,fill=blue!5,rounded corners=2pt,text width=3.3cm,minimum height=1.5cm,align=center,font=\small,inner sep=5pt},
  flow/.style={-{Stealth[length=5pt]},thick,draw=black!60},
  lab/.style={font=\scriptsize,text=black!60,fill=white,inner sep=2pt}]
\node[box] (v) at (1.7,0) {\textbf{已核验输入}\\带日期的数据快照};
\node[box] (g) at (6.6,0) {\textbf{执行候选方案}\\保存的输出};
\node[box,draw=forestgreen,fill=forestgreen!6] (c) at (11.5,0) {\textbf{审查证据}\\验收或修复};
\draw[flow] (v)--node[above,lab]{就绪}(g);
\draw[flow] (g)--node[above,lab]{产物}(c);
\draw[flow,draw=darkred] (c.south)--(11.5,-1.7)--node[below,lab,text=darkred]{审查未通过：交回问题项}(6.6,-1.7)--(g.south);
\node[font=\small,text=black!60] at (6.6,-2.8) {前置条件未满足：先等待，再分派任务。};
\end{tikzpicture}
}
\end{center}
\begin{itemize}
\item 节点执行任务；边规定前置条件与交接方式。
\item 缺少输入时先等待；审查未通过时，将结果交回 Generator。
\item 前向任务链可以无环；加入修复边后，完整工作流便包含循环。
\end{itemize}
\end{frame}

% Page 26
\begin{frame}{图：产物与共享状态}
\begin{center}
\resizebox{.96\textwidth}{!}{\begin{tikzpicture}[
  box/.style={draw=blue,fill=blue!5,rounded corners=2pt,align=center,font=\small,inner sep=6pt},
  flow/.style={-{Stealth[length=5pt]},thick,draw=black!60}]
\node[box,text width=4.2cm,minimum height=3.1cm] (w) at (2.3,0) {\textbf{工作 Agent 产物}\\[4pt]
请求 ID\\
候选方案 + 版本\\
状态 + 问题项\\
证据路径 + 哈希值};
\node[box,text width=3.3cm,minimum height=3.1cm,draw=forestgreen,fill=forestgreen!6] (v) at (7.0,0) {\textbf{交接核验}\\[4pt]
结构约束 + 身份\\
完整报告\\
证据完整性};
\node[box,text width=3.5cm,minimum height=3.1cm] (s) at (11.4,0) {\textbf{控制器状态}\\[4pt]
待处理角色 + 轮次\\
结转的未解决问题项\\
下一项满足条件的任务};
\draw[flow] (w)--(v);
\draw[flow] (v)--node[above,font=\scriptsize]{通过}(s);
\draw[flow,draw=darkred] (v.south)--(7.0,-2.3)--node[below,font=\scriptsize]{无效：修复报告；不得推进}(2.3,-2.3)--(w.south);
\end{tikzpicture}
}
\end{center}
\begin{itemize}
\item 各工作 Agent 分别编写报告。只有控制器可以更新运行的控制状态。
\item 合并前核验身份、完整性与证据。
\item 明确保留相互冲突的问题项；最后提交的报告不得悄然将其覆盖。
\end{itemize}
\end{frame}

% Page 27
\begin{frame}{图：团队、任务分发与汇合}
\setlength{\topsep}{0pt}
\begin{center}
\resizebox{.86\textwidth}{!}{\begin{tikzpicture}[
  lead/.style={draw,rounded corners=2pt,line width=.7pt,align=center,font=\small,inner sep=4pt},
  worker/.style={draw,dashed,rounded corners=2pt,fill=white,text width=1.30cm,minimum height=.65cm,align=center,font=\footnotesize,inner sep=2pt},
  flow/.style={-{Stealth[length=4pt]},line width=.65pt,draw=black!60},
  label/.style={font=\footnotesize,align=center,text=black!60,fill=white,inner sep=1pt}]
\node[lead,draw=forestgreen,fill=forestgreen!6,text width=5.0cm] (o) at (7.65,0)
  {\textbf{Orchestrator}\\状态 + 检查关卡 + 路由};
\node[font=\footnotesize,align=center] (idea) at (1.3,0) {\textbf{用户想法}};
\draw[flow] (idea.east)--(o.west);

\draw[blue,line width=.7pt,rounded corners=3pt,fill=blue!4] (0,-1.1) rectangle (5.4,-3.95);
\draw[darkred,line width=.7pt,rounded corners=3pt,fill=red!4] (7.8,-1.1) rectangle (15.3,-3.95);
\node[font=\small\bfseries,text=blue] at (2.7,-1.38) {Generator 负责人};
\node[font=\small\bfseries,text=darkred] at (11.55,-1.38) {Critic 负责人};
\node[font=\footnotesize,text=black!60] at (2.7,-1.72) {协调构建专项 Agent};
\node[font=\footnotesize,text=black!60] at (11.55,-1.72) {每轮分派新上下文的审查 Agent};

\node[worker,draw=blue] (gd) at (1.0,-2.48) {数据};
\node[worker,draw=blue] (gf) at (2.7,-2.48) {特征};
\node[worker,draw=blue] (gb) at (4.4,-2.48) {回测};
\node[worker,draw=darkred,text width=1.55cm] (cd) at (8.8,-2.48) {数据与证据};
\node[worker,draw=darkred,text width=1.55cm] (cr) at (10.65,-2.48) {稳健性};
\node[worker,draw=darkred,text width=1.55cm] (ce) at (12.5,-2.48) {执行};
\node[worker,draw=darkred,text width=1.55cm] (ck) at (14.35,-2.48) {风险与控制};
\node[lead,draw=blue,fill=white,text width=4.3cm,minimum height=.5cm] (gi) at (2.7,-3.5)
  {\textbf{整合 + 测试}};
\node[lead,draw=darkred,fill=white,text width=6.4cm,minimum height=.5cm] (cm) at (11.55,-3.5)
  {\textbf{合并 + 覆盖检查}};
\foreach \worker in {gd,gf,gb}{\draw[flow] (2.7,-1.96)--(\worker.north);\draw[flow] (\worker.south)--(gi.north);}
\foreach \worker in {cd,cr,ce,ck}{\draw[flow] (11.55,-1.96)--(\worker.north);\draw[flow] (\worker.south)--(cm.north);}

\draw[flow] (o.south west)--(4.7,-.75)--(2.7,-.75)--(2.7,-1.1);
\draw[flow] (o.south east)--(10.3,-.75)--(11.55,-.75)--(11.55,-1.1);
\node[label] at (2.7,-.55) {分派};
\node[label] at (11.55,-.55) {分派};
\draw[flow,draw=blue] (5.4,-2.0)--(7.8,-2.0);
\node[label,text=blue] at (6.6,-1.62) {核验后的输入\\经由 Orchestrator};
\draw[flow,dashed,draw=darkred] (7.8,-2.95)--(5.4,-2.95);
\node[circle,draw=darkred,fill=white,inner sep=1pt,text=darkred,font=\small] at (6.6,-2.95) {$\times$};
\node[label,text=darkred] at (6.6,-3.37) {审查结果经由\\Orchestrator};

\node[lead,draw=forestgreen,fill=forestgreen!6,text width=5.2cm] (packet) at (7.65,-4.75)
  {\textbf{Orchestrator：修复任务包}\\台账 + 决策：继续 / 结束};
\draw[flow] (cm.south)--(11.55,-4.75)--(packet.east);
\node[label,anchor=west] at (11.7,-4.35) {合并后的问题项};
\draw[flow] (packet.west)--(-.65,-4.75)--(-.65,-1.38)--(0,-1.38);
\node[label] at (2.7,-4.35) {继续时下发修复任务包};
\node[font=\footnotesize,align=left,text=black!60,anchor=west] at (10.4,-5.2) {实线：角色负责人\\虚线：专项 Agent};
\end{tikzpicture}
}
\end{center}
\begin{itemize}
\item 检查需要不同专业知识时，使用专项 Agent。
\item 等待所有必需报告提交，再核验身份与证据，之后才能合并结果。
\end{itemize}
\end{frame}

% Page 28
\begin{frame}{图：合并结果与共同限额}
\begin{center}
\resizebox{.93\textwidth}{!}{\begin{tikzpicture}[box/.style={draw=blue,fill=blue!5,rounded corners=2pt,text width=3.6cm,minimum height=1.3cm,align=center,font=\small},link/.style={-{Stealth[length=5pt]},thick,draw=black!60}]
\node[box] (a) at (2,0) {\textbf{方案 A}\\满足各自的局部限额};
\node[box] (b) at (2,-2.1) {\textbf{方案 B}\\满足各自的局部限额};
\node[box,draw=darkred,fill=red!4] (total) at (10,-1.05) {\textbf{合并风险敞口}\\可能突破共同上限};
\draw[link] (a.east)--(total.north west);
\draw[link] (b.east)--(total.south west);
\end{tikzpicture}
}
\end{center}
\begin{center}\illustrative\end{center}
\begin{itemize}
\item 两项方案可各自满足局部限额，但合并后仍超出投资组合的共同限额。
\item 课堂测试样例尚未实施这项组合层面的控制，需要另行补充。
\end{itemize}
\end{frame}

% Page 29
\begin{frame}{图：条件路由与恢复}
\begin{center}
\small
\begin{tabular}{@{}P{3.6cm}P{4.2cm}P{4.7cm}@{}}
\toprule
\textbf{已核验状态} & \textbf{路由路径} & \textbf{后续依赖任务} \\
\midrule
输入缺失 & 等待前置条件 & 暂不分派任务 \\
\rowcolor{gray!10}
提交无效 & 由当前角色修复 & 继续阻止下游推进 \\
审查要求修复 & 将问题项交回 Generator & 重新核验修订后输出 \\
\rowcolor{gray!10}
必要检查通过 & 转向下一项满足条件的任务 & 汇合全部必需报告 \\
终止结果 & 结束并保留原因 & 不再分派任务 \\
\rowcolor{gray!10}
获准的输入变更 & 启动新案例 & 重新生成依赖这些输入的输出 \\
\bottomrule
\end{tabular}

\end{center}
\begin{itemize}
\item 依据已核验的状态与依赖关系，选择下一项任务。
\item 输入版本发生变化时启动新案例：保留旧证据，并重新构建依赖这些输入的工作。
\end{itemize}
\end{frame}

\subsection{演进}

% Page 30
\begin{frame}{演进：哪些内容可以改变？}
\begin{center}
\small
\begin{tabular}{@{}P{3.3cm}P{5.3cm}P{3.9cm}@{}}
\toprule
\textbf{改动对象} & \textbf{示例} & \textbf{所需证据} \\
\midrule
提示词 / 上下文 & 要求说明证据缺失 & 行为比较 \\
\rowcolor{gray!10}
技能 & 增加来源交叉核验步骤 & 流程测试 \\
工具 / 代码 & 返回明确的失败状态 & 单元测试与集成测试 \\
\rowcolor{gray!10}
角色 / 工作流 & 增加独立审查路径 & 端到端评估 \\
\bottomrule
\end{tabular}

\end{center}
\begin{itemize}
\item 工具循环选择行动；研究循环修订候选策略。
\item 演进评估可复用的系统改动，并决定采用、否决或暂缓。
\item 策略参数搜索仍属于研究过程。QT 通过 Generator 提示词的版本更新演示演进。
\end{itemize}
\end{frame}

% Page 31
\begin{frame}{演进：改进系统的完整流程}
\setlength{\topsep}{0pt}
\begin{center}
\resizebox{.96\textwidth}{!}{\begin{tikzpicture}[
  box/.style={draw=blue,fill=blue!5,rounded corners=2pt,text width=2.65cm,minimum height=1.4cm,align=center,font=\footnotesize,inner sep=4pt},
  flow/.style={-{Stealth[length=4pt]},line width=.8pt,draw=black!60}]
\node[draw=forestgreen,fill=forestgreen!5,rounded corners=2pt,font=\small,align=center,text width=12.9cm,inner sep=5pt] at (6.75,1.2)
{\textbf{演进 Orchestrator：}失败、改动范围、变更前版本、预算与状态};
\node[box] (e) at (1.5,0) {\textbf{已完成案例}\\保留的证据};
\node[box] (d) at (5.0,0) {\textbf{诊断}\\反复出现的失败};
\node[box] (g) at (8.5,0) {\textbf{Generator 团队}\\限定范围的改动\\及测试};
\node[box] (t) at (12.0,0) {\textbf{评估}\\固定案例 + 回归检查};
\node[box,draw=darkred,fill=red!4] (c) at (12.0,-1.85) {\textbf{Critic 团队}\\新上下文\\证据、范围、回归检查};
\node[box,draw=darkred,fill=red!4] (h) at (8.5,-1.85) {\textbf{系统负责人决策}\\批准 / 否决\\或暂缓};
\node[box,draw=forestgreen,fill=forestgreen!6] (v) at (5.0,-1.85) {\textbf{版本注册表}\\采用或保留先前版本};
\node[box,draw=forestgreen,fill=forestgreen!6] (n) at (1.5,-1.85) {\textbf{后续案例}\\当前启用版本};
\draw[flow] (e)--(d);\draw[flow] (d)--(g);\draw[flow] (g)--(t);
\draw[flow] (t)--node[right,font=\scriptsize]{通过}(c);
\draw[flow] (c)--node[above,font=\scriptsize]{通过}(h);
\draw[flow] (h)--(v);\draw[flow] (v)--(n);
\draw[flow,dashed] (n.west)--(-.35,-1.85)--(-.35,0)--(e.west);
\node[font=\scriptsize,text=darkred,align=center] at (6.75,-2.9) {检查或审查未通过：保留当前启用版本及证据。仍可回滚。};
\end{tikzpicture}
}
\end{center}
\begin{itemize}
\item Generator 提出限定范围的改动与测试；新上下文中的 Critic 团队审查证据、回归检查与改动范围。
\item 只有通过必要检查与审查的版本才可获批。改动若被否决或暂缓，继续保留当前启用版本。
\end{itemize}
\end{frame}

% Page 32
\begin{frame}{演进：证据与评估}
\begin{center}
\small
\begin{tabular}{@{}P{3.4cm}P{9.1cm}@{}}
\toprule
\textbf{评估标准} & \textbf{要求的行为} \\
\midrule
格式 & 必需字段能够正确解析。 \\
\rowcolor{gray!10}
语义 & 行为满足任务的实际约束。 \\
证据 & 引用的来源存在于保留证据中，且能够支持相关论断。 \\
\rowcolor{gray!10}
授权 & 未经授权的操作仍被阻止。 \\
\bottomrule
\end{tabular}

\end{center}
\begin{center}
\resizebox{.93\textwidth}{!}{\begin{tikzpicture}[
 box/.style={draw=blue,fill=blue!5,rounded corners=2pt,text width=5.6cm,minimum height=1.0cm,align=center,font=\small,inner sep=5pt}]
\node[box] at (3.0,0) {\textbf{开发案例}\\反馈可用于修订};
\node[box,draw=forestgreen,fill=forestgreen!6] at (10.0,0) {\textbf{留出案例}\\用于独立的最终评估};
\end{tikzpicture}
}
\end{center}
\begin{itemize}
\item 比较改动之前，先确定评估标准。
\item QT 的时序检查要求执行晚于信号实际可用时点；仅有日期字段并不充分。
\end{itemize}
\end{frame}

% Page 33
\begin{frame}{演进：回归检查与独立审查}
\begin{center}
\small
\begin{tabular}{@{}P{4.0cm}P{8.5cm}@{}}
\toprule
\textbf{审查对象} & \textbf{所需证据} \\
\midrule
控制器行为 & 报告缺失、过期 ID 或无效证据仍能阻止流程推进。 \\
\rowcolor{gray!10}
模型行为 & 根据时序与证据检查，比较新旧输出。 \\
改动范围 & 改动对象及权限与获批的改动请求一致。 \\
\rowcolor{gray!10}
独立审查 & 核验改动已修复记录的失败，且在预留案例中保持既有行为。 \\
\bottomrule
\end{tabular}

\end{center}
\begin{itemize}
\item 代码测试不能证明模型行为得到改善。
\item QT 的提示词接口不接受对风险限额、策略参数和控制器代码的修改。
\item 更广泛的改动需要另行界定范围与测试；改动提案不能降低自身的验收标准。
\end{itemize}
\end{frame}

% Page 34
\begin{frame}{演进：版本、批准与回滚}
\begin{center}
\resizebox{.94\textwidth}{!}{\begin{tikzpicture}[
 box/.style={draw=blue,fill=blue!5,rounded corners=2pt,text width=3.4cm,minimum height=1.15cm,align=center,font=\small,inner sep=5pt},
 flow/.style={-{Stealth[length=4pt]},line width=.8pt,draw=black!60},
 lab/.style={font=\scriptsize,text=black!60,fill=white,inner sep=1pt,align=center}]
\node[box] (old) at (1.8,0) {\textbf{当前版本}\\保留的配置};
\node[box] (cand) at (6.8,0) {\textbf{候选版本}\\差异 + 测试证据};
\node[box,draw=forestgreen,fill=forestgreen!6] (active) at (11.8,0) {\textbf{获批版本}\\新的启用版本指针};
\draw[flow] (old)--node[above,lab]{提议\\+ 测试}(cand);
\draw[flow] (cand)--node[above,lab]{审查\\+ 批准}(active);
\node[box,fill=white,minimum height=.85cm] (existing) at (1.8,-1.8) {\textbf{现有案例}\\保留各自的\\原始版本};
\node[box,fill=white,minimum height=.85cm] (new) at (11.8,-1.8) {\textbf{新案例}\\使用当前启用版本};
\draw[flow] (old)--(existing);\draw[flow] (active)--(new);
\draw[flow,draw=darkred,dashed] (active.north)--(11.8,1.25)--node[above,lab,text=darkred]{获准回滚：恢复先前的版本指针}(1.8,1.25)--(old.north);
\end{tikzpicture}
}
\end{center}
\begin{itemize}
\item 为每个版本记录测试时使用的模型、工具、数据、规则和提示词。
\item 别名用于选择版本；批准仍需独立记录。
\item 保留新旧两个版本及其证据。回滚会恢复先前的版本指针，并须在新案例中核验行为。
\end{itemize}
\end{frame}

% Page 35
\begin{frame}{演进：监测与改动触发条件}
\begin{center}
\resizebox{.94\textwidth}{!}{\begin{tikzpicture}[
  box/.style={draw=blue,fill=blue!5,rounded corners=2pt,text width=3.3cm,minimum height=1.5cm,align=center,font=\small,inner sep=5pt},
  link/.style={-{Stealth[length=5pt]},thick,draw=black!60}]
\node[box] (n0) at (1.7,0) {后续案例\\形成新证据};
\node[box] (n1) at (6.6,0) {监测故障\\与共同风险};
\draw[link] (n0) -- (n1);
\node[box,draw=forestgreen,fill=forestgreen!6] (n2) at (11.5,0) {审查变更\\再决定采纳};
\draw[link] (n1) -- (n2);
\end{tikzpicture}
}
\end{center}
\begin{itemize}
\item 监测反复出现的失败与共同风险。
\item 模型、工具、数据源或权限变化时，重新评估。
\item 将已记录的失败反馈给演进流程。提示词优化并不保证可靠性。
\end{itemize}
\end{frame}

% Page 36
\begin{frame}[t]{完整架构：各部分如何协同}
\centering
\resizebox{!}{6.0cm}{\begin{tikzpicture}[
  lead/.style={draw,rounded corners=2pt,line width=.7pt,align=center,font=\footnotesize,inner sep=4pt},
  worker/.style={draw,dashed,rounded corners=2pt,fill=white,text width=1.30cm,minimum height=.62cm,align=center,font=\scriptsize,inner sep=2pt},
  flow/.style={-{Stealth[length=4pt]},line width=.65pt,draw=black!60},
  edge/.style={font=\scriptsize,align=center,text=black!65,fill=white,inner sep=1pt}]
\draw[darkorange!85!black,line width=1.1pt,rounded corners=5pt]
  (-.45,.45) rectangle (15.75,-5.65);
\node[fill=white,inner sep=4pt,font=\small\bfseries,text=darkorange!85!black]
  at (7.65,.45) {演进：经评估的系统改动};
\node[lead,draw=forestgreen,fill=forestgreen!6,text width=4.5cm,minimum height=.68cm] (o) at (7.65,-.20)
  {\textbf{Orchestrator}\\分派 + 状态 + 路由};
\node[font=\footnotesize,align=center,text width=3.1cm] (idea) at (1.5,.85)
  {\textbf{用户想法}\\任务授权 + 限额};
\node[font=\scriptsize,align=center,text=darkred,text width=3.4cm] (human) at (13.6,.85)
  {\textbf{人工制定规则}\\并保留批准权};
\draw[flow] (idea.south)--(1.5,-.20)--(o.west);
\draw[flow,draw=darkred] (human.south)--(13.6,-.20)--(o.east);

\draw[blue,line width=.7pt,rounded corners=3pt,fill=blue!4] (0,-1.0) rectangle (5.4,-3.7);
\draw[darkred,line width=.7pt,rounded corners=3pt,fill=red!4] (7.8,-1.0) rectangle (15.3,-3.7);
\node[font=\footnotesize\bfseries,text=blue] at (2.7,-1.28) {Generator 负责人};
\node[font=\footnotesize\bfseries,text=darkred] at (11.55,-1.28) {Critic 负责人};
\node[font=\scriptsize,text=black!65] at (2.7,-1.60) {协调构建专项 Agent};
\node[font=\scriptsize,text=black!65] at (11.55,-1.60) {每轮开展全新审查};
\node[worker,draw=blue] (gd) at (1.0,-2.30) {数据};
\node[worker,draw=blue] (gf) at (2.7,-2.30) {特征};
\node[worker,draw=blue] (gb) at (4.4,-2.30) {回测};
\node[worker,draw=darkred,text width=1.52cm] (cd) at (8.8,-2.30) {数据与证据};
\node[worker,draw=darkred,text width=1.52cm] (cr) at (10.65,-2.30) {稳健性};
\node[worker,draw=darkred,text width=1.52cm] (ce) at (12.5,-2.30) {执行};
\node[worker,draw=darkred,text width=1.52cm] (ck) at (14.35,-2.30) {风险与控制};
\node[lead,draw=blue,fill=white,text width=4.3cm,minimum height=.52cm] (gi) at (2.7,-3.18)
  {\textbf{整合 + 测试}};
\node[lead,draw=darkred,fill=white,text width=6.4cm,minimum height=.52cm] (cm) at (11.55,-3.18)
  {\textbf{合并 + 覆盖检查}};
\foreach \w in {gd,gf,gb}{\draw[flow] (2.7,-1.83)--(\w.north);\draw[flow] (\w.south)--(gi.north);}
\foreach \w in {cd,cr,ce,ck}{\draw[flow] (11.55,-1.83)--(\w.north);\draw[flow] (\w.south)--(cm.north);}
\draw[flow] (o.south west)--(4.7,-.82)--(2.7,-.82)--(2.7,-1.0);
\draw[flow] (o.south east)--(10.3,-.82)--(11.55,-.82)--(11.55,-1.0);
\node[edge,anchor=south] at (2.7,-.75) {分派};
\node[edge,anchor=south] at (11.55,-.75) {检查通过后};
\draw[flow,draw=blue] (5.4,-2.0)--(7.8,-2.0);
\node[edge,text=blue] at (6.6,-1.62) {可执行检查关卡\\经由 Orchestrator};
\node[lead,draw=forestgreen,fill=forestgreen!6,text width=4.6cm,minimum height=.65cm] (packet) at (7.65,-4.30)
  {\textbf{Orchestrator：修复任务包}\\问题项 + 决策：继续 / 结束};
\draw[flow] (cm.south)--(11.55,-4.30)--(packet.east);
\node[edge,anchor=west] at (11.7,-4.1) {问题项交回\\Orchestrator};
\draw[flow] (packet.west)--(2.7,-4.30)--(2.7,-3.7);
\node[edge,anchor=east] at (2.55,-4.1) {继续时\\发回修复};
\node[draw=black!35,rounded corners=2pt,fill=gray!5,text width=14.4cm,align=center,
  minimum height=.45cm,font=\scriptsize,inner sep=4pt] at (7.65,-5.15)
  {\textbf{共享状态与证据：}请求 ID、产物、哈希值、问题项与版本};
\end{tikzpicture}
}
\par\raggedright
\smallskip
\begin{itemize}
\item 在图中找出各角色的 Harness、修订循环、交接路径与演进边界。
\end{itemize}
\end{frame}
